Este capítulo recoge la arquitectura física planteada para el sistema, las instrucciones para su despliegue y las instrucciones para la operación y mantenimiento del nivel de servicio.

\section{Arquitectura Física}
En este apartado, describimos los principales elementos \textit{hardware} que forman la arquitectura física de nuestro sistema. Se debe incluir un modelo de despliegue en el cual se describe cómo los elementos \textit{software} son instalados y ejecutados en los elementos \textit{hardware} (físicos o virtuales). También se incluyen las especificaciones y los requisitos del \textit{hardware} (servidores, etc.), así como de los elementos \textit{software} (sistemas operativos, servicios, aplicaciones, etc.) necesarios.

\section{Instrucciones de despliegue}
A continuación, se describen las instrucciones de instalación del sistema sobre la infraestructura física descrita anteriormente.

\subsection{Requisitos previos}
Requisitos de hardware y \textit{software} para la correcta instalación del sistema.

\subsection{Inventario de componentes}
Lista de los componentes hardware y software que se incluyen en la versión del producto.

\subsection{Procedimientos de instalación}
Procedimientos de instalación y configuración de cada componente hardware y software (base y desarrollado) para asegurar la correcta instalación y explotación del sistema, así como aquellos procedimientos necesarios de migración/carga de datos.

\subsection{Pruebas de implantación}
Descripción de las pruebas a realizar después de la instalación del sistema. 

\section{Instrucciones para la operación del sistema y mantenimiento del nivel de servicio}
Procedimientos necesarios para asegurar el correcto funcionamiento, rendimiento, disponibilidad y seguridad del sistema: back-ups, chequeo de logs, etc. También es preciso indicar claramente aquellas actuaciones precisas necesarias para el mantenimiento preventivo del sistema y así prevenir posibles fallos en el mismo. 